% -------------------- ŞEKİLLER DİZİNİ --------------------
\chapter*{ŞEKİLLER DİZİNİ}
\addcontentsline{toc}{chapter}{ŞEKİLLER DİZİNİ}

\begin{longtable}{p{2cm}p{11cm}r}
\textbf{Şekil} & \textbf{Açıklama} & \textbf{Sayfa} \\ \hline
\hyperref[fig:sistem_mimarisi]{Şekil 1.1} & DeepSeis Sistem Mimarisi Genel Görünümü & \pageref{fig:sistem_mimarisi} \\
\hyperref[fig:sismik_dalgalar]{Şekil 2.1} & Sismik Dalga Tipleri ve Yayılım Karakteristikleri & \pageref{fig:sismik_dalgalar} \\
\hyperref[fig:ml_anomaly_pipeline]{Şekil 2.2} & STFT Öznitelikleriyle Sınıflandırma Tabanlı Anomali Tespiti Yaklaşımı & \pageref{fig:ml_anomaly_pipeline} \\
\hyperref[fig:stft_ornek]{Şekil 3.1} & STFT Tabanlı Zaman-Frekans Dönüşümü Örneği & \pageref{fig:stft_ornek} \\
\hyperref[fig:clean_model_comparison]{Şekil 4.1} & Temiz Normal/Anomali Ayrımı Model Karşılaştırması & \pageref{fig:clean_model_comparison} \\
\hyperref[fig:clean_confusion]{Şekil 4.2} & HistGradientBoosting Modeli Karışıklık Matrisi & \pageref{fig:clean_confusion} \\
\hyperref[fig:clean_roc]{Şekil 4.3} & HistGradientBoosting Modeli ROC Eğrisi & \pageref{fig:clean_roc} \\
\hyperref[fig:clean_pr]{Şekil 4.4} & HistGradientBoosting Modeli Precision-Recall Eğrisi & \pageref{fig:clean_pr} \\
\end{longtable}
